\chapter{Ensembles, dénombrement, dénombrabilité}

\section{Vocabulaire de la logique}

$\exists x$: il existe $x$.

$\forall x$: pour tout $x$.

Pour deux propositions $P$ et $Q$:
\begin{itemize}
    \item $P \land Q$: $P$ et $Q$.
    \item $P \lor Q$: $P$ ou $Q$.
    \item $\neg P$: non $P$ / la négation de $P$.
    \item $P \implies Q$: $P$ implique $Q$;
        Si $P$, alors $Q$;
        $P$ est une condition suffisante de $Q$;
        $Q$ est une condition nécessaire de $P$.
    \item $P \iff Q$: $P$ et $Q$ sont équivalentes;
        $P$ si et seulement si (ssi) $Q$;
        $P$ est une condition nécessaire et suffisante (CNS) de $Q$.
\end{itemize}

\begin{itemize}
    \item $Q \implies P$: la réciproque de $P \implies Q$.
    \item $\neg P \implies \neg Q$: l'inverse de $P \implies Q$.
    \item $\neg Q \implies \neg P$: la contraposée de $P \implies Q$.
\end{itemize}

Notations (de Bourbaki) :
$\bbN$: L'ensemble des entiers naturels.
$\bbN = \{ 0, 1, 2, \ldots \}, \bbN^* = \bbN \backslash \{ 0 \}$

$\bbZ$: L'ensemble des entiers relatifs.
$\bbZ = \{ \ldots, -2, -1, 0, 1, 2, \ldots \}$

$\bbQ$: L'ensemble des nombres rationnels.
$\bbQ = \{ \frac{a}{b} : a \in \bbZ, b \in \bbZ^* \}$

$\bbR$: L'ensemble des nombres réels.

$\bbC$: L'ensemble des nombres complexes.

\section{Opérations sur les ensembles}

Un ensemble = Une collection d'objets (définition ``naïve'' / non axiomatique)

Deux ensembles sont égaux ssi ils contiennent les mêmes éléments.

L'ensemble vide est noté $\varnothing$.

Soit $E$ un ensemble :
\begin{itemize}
    \item $x \in E$: $x$ appartient à $E$; $x$ est un élément de $E$.
    \item $A \subset E / A \subseteq E$: $\forall w \in A$, on a $w \in E$;
        $A$ est inclus dans $E$;
        $A$ est un sous-ensemble de $E$;
        $A$ est une partie de $E$.
\end{itemize}

Soit $A, B$ deux parties de $E$:
\begin{itemize}
    \item $A \cup B := \{ w \in E: w \in A \text{ ou } w \in B \}$. union / réunion de $A$ et $B$.
    \item $A \cap B := \{ w \in E: w \in A \text{ et } w \in B \}$. intersection de $A$ et $B$.
    \item $A^C := \{ w \in E: w \notin A \}$. complémentaire de $A$ (par rapport à $E$).
    \item $A \backslash B := A \cap (B^C)$. différence de $A$ et $B$.
    \item $A \triangle B := (A \backslash B) \cup (B \backslash A)$. différence symétrique de $A$ et $B$.
\end{itemize}

Soit $I$ un ensemble d'indices, pour tout $i \in I$, soit $A_i$ un sous-ensemble de $E$:
\begin{itemize}
    \item $(A_i)_{i \in I}$: famille de sous-ensembles de $E$.
    \item $\bigcup\limits_{i \in I} A_i := \{ w \in E: \exists i \in I \text{ tel que } w \in A_i \}$.
        réunion de la famille $(A_i)_{i \in I}$.
    \item $\bigcap\limits_{i \in I} A_i := \{ w \in E: \forall i \in I, w \in A_i \}$.
        intersection de la famille $(A_i)_{i \in I}$.
\end{itemize}

La famille $(A_i)_{i \in I}$ est \textcolor{red}{disjointe} si
\[ \forall i, j \in I \text{ tels que } i \neq j, on a A_i \cap A_j = \varnothing \]

Distributivité entre $\cup, \cap, (\quad)^C$:
\begin{itemize}
    \item $\left( \bigcup\limits_{i \in I} A_i \right) \cap B = \bigcup\limits_{i \in I} (A_i \cap B)$
    \item $\left( \bigcap\limits_{i \in I} A_i \right) \cup B = \bigcap\limits_{i \in I} (A_i \cup B)$
    \item $\left( \bigcup\limits_{i \in I} A_i \right)^C = \bigcap\limits_{i \in I} (A_i^C)$
    \item $\left( \bigcap\limits_{i \in I} A_i \right)^C = \bigcup\limits_{i \in I} (A_i^C)$
\end{itemize}

Produit cartésien :
\[ A \times B := \{ (a, b): a \in A, b \in B \} \]

Pour $n \geq 2$ entier, on écrit $A^n = \underbrace{A \times A \times A \times \cdots \times A}_{n \text{ fois}}$.

$\calP(E) := \{ A: A \subseteq E \}$: l'ensemble des sous-ensembles de $E$.

\begin{rqq}
    On a toujours $\varnothing \in \calP(E)$ et $E \in \calP(E)$.
    $\varnothing$ et $E$ sont appelés sous-ensembles triviaux de $E$.
\end{rqq}

$\forall A \in \calP(E)$, on définit la fonction indicatrice de $A$ par
\[ \bbI_A: E \to \bbR, \bbI_A(x) = \begin{cases} 1, & \text{si} x \in A \\ 0, & \text{si} x \notin A \end{cases} \]

\begin{deff}
    Une fonction $f: E \to F$ est
    \begin{itemize}
        \item \textbf{injective}, si tout élément de $F$ est l'image d'\textcolor{red}{au plus} un élément de $E$.
            càd $\forall y \in F$, il existe au plus un élément $x \in E$ tel que $y = f(x)$.
        \item \textbf{surjective}, si tout élément de $F$ est l'image d'\textcolor{red}{au moins} un élément de $E$.
            càd $\forall y \in F, \exists x \in E$ tel que $y = f(x)$.
        \item \textbf{bijective}, si $f$ est à la fois injective et surjective.
    \end{itemize}
\end{deff}

Si $f: E \to F$ est bijective, alors on peut définir son inverse $f^{-1}: F \to E$.

Soit $B \subseteq F$, l'image réciproque de $B$ par $f$ est $f^{-1}(B)$.

Soit $A \subseteq E$, l'image directe de $A$ par $f$ est $f(A)$.

\subsection{Ensembles finis}

Notation: Pour $n \in \bbN$, on note $\{1, 2, \ldots, n\} = \llb 1, n \rrb$ et
par convention si $n = 0, \llb 1, 0 \rrb = \varnothing$.

\begin{deff}
    Un ensemble $E$ est fini s'il existe $n \in \bbN$ et une bijection $f: \llb 1, n \rrb \to E$.
\end{deff}

\begin{lem}[fondamental]
    Soient $m, n \in \bbN$, s'il existe une injection $f: \llb 1, n \rrb \to \llb 1, m \rrb$, alors $n \leq m$.
\end{lem}

\begin{proof}
    On raisonne par récurrence sur $m$.
    \begin{itemize}
        \item (Initialisation)
            Si $m = 0$, alors $\llb 1, m \rrb$ est vide.
            L'application $f: \llb 1, n \rrb \to \varnothing$ peut exister seulement si
            $\llb 1, n \rrb \neq \varnothing$.
            Donc $n \leq m$ est vrai.
        \item (Hérédité)
            Supposons que le résultat est vrai pour $m$.
            Soit $f: \llb 1, n \rrb \to \llb 1, m + 1 \rrb$ une injection.
            Si $n = 0$, alors $n \leq m + 1$ et c'est fini.
            Supposons $n \geq 1$.
            Soit $p: \llb 1, m + 1 \rrb \to \llb 1, m + 1 \rrb$ la permutation qui échange $f(n)$ et $m + 1$.
            Alors $g = p \circ f$ est une application $g: \llb 1, n \rrb \to \llb 1, m + 1 \rrb$ telle que $g(n) = m + 1$.
            De plus, $g$ est injective.
            La restriction de $g$ sur $\llb 1, n - 1 \rrb$ est une injection $\llb 1, n - 1 \rrb \to \llb 1, m \rrb$.
            Par l'hypothèse de récurrence, on a $n - 1 \leq m$, donc $n \leq m + 1$.
    \end{itemize}

    Donc le lemme est démontré par récurrence. \qedhere
\end{proof}

\begin{deff}
    Pour un ensemble fini $E$.
    S'il existe une bijection $f: \llb 1, n \rrb \to E$ pour l'entier $n$,
    alors $n$ est appelé le cardinal de $E$, et on écrit $card(E) = |E| = \#E = n$.
\end{deff}

\begin{lem}
    Soit $E$ un ensemble fini. L'entier $n$ tel qu'il existe une bijection $\llb 1, n \rrb \to E$ est unique.
\end{lem}

\begin{proof}
    Supposons qu'il existe deux bijections
    \[ f: \llb 1, n \rrb \to E \quad \text{et} \quad g: \llb 1, m \rrb \to E \]
    Alors $g^{-1} \circ f$ est une bijection de $\llb 1, n \rrb$ dans $\llb 1, m \rrb$.
    En particulier, c'est une injection $\llb 1, n \rrb \to \llb 1, m \rrb$.
    Selon le lemme fondamental, on a $n \leq m$.
    De la même façon, $m \leq n$.
    Donc $n = m$, d'où l'unicité de la valeur de $n$. \qedhere
\end{proof}

Grâce au lemme précédent, la définition du cardinal pour les ensembles finis est bien posée.

\begin{rqq}
    L'ensemble vide $\varnothing$ est le seul ensemble de cardinal nul.
\end{rqq}

Si $card(E) = n$, alors on peut représenter les éléments de $E$ comme une série ordonnée finie
$E = \{ a_1, a_2, \ldots, a_n \}$ avec $a_i \neq a_j, \forall i \neq j$.

\begin{deff}
    Un ensemble infini est un ensemble qui n'est pas fini.
\end{deff}

\begin{exem}
    $\bbN, \bbZ, \bbQ, \bbR, \bbC$ sont tous des ensembles infinis.
\end{exem}

\begin{thm}[comparaison]
    Soient $E, F$ deux ensembles finis.
    \begin{enumerate}
        \item $|E| \leq |F|$ ssi $\exists$ une injection $f: E \to F$
        \item $|E| = |F|$ ssi $\exists$ une bijection $f: E \to F$
        \item Si $E \neq \varnothing$, $|E| \leq |F|$ ssi $\exists$ une surjection $F \to E$.
    \end{enumerate}
\end{thm}

\begin{proof}
    Soient $n = |E|, m = |F|$ et $u: \llb 1, n \rrb \to E, v: \llb 1, m \rrb \to F$ deux bijections.
    \begin{enumerate}
        \item
            Supposons qu'il existe une injection $f: E \to F$.
            Alors l'application $v^{-1} \circ f \circ u: \llb 1, n \rrb \to \llb 1, m \rrb$ est injective.
            Selon le lemme fondamental, on a $n \leq m$, càd $|E| \leq |F|$.
            Réciproquement, si $n \leq m$, alors l'application $h: \begin{array}[t]{ccc}
            \llb 1, n \rrb & \to & \llb 1, m \rrb \\ k & \mapsto & k
            \end{array} $ est injective.
            Donc $v \circ h \circ u^{-1}: E \to F$ est une injection.
            Donc il existe une injection $E \to F$.
        \item
            Supposons qu'il existe une bijection $f: E \to F$.
            Puisque $f$ est injective, selon le point 1, on a $|E| \leq |F|$.
            De même, $f^{-1}: F \to E$ est aussi injective.
            Donc $|F| \leq |E|$.
            Donc $|E| = |F|$.
        \item
            Supposons $E \neq \varnothing$, càd $n \geq 1$.
            Supposons qu'il existe une surjection $g: F \to E$.
            Alors $q = v \circ g \circ u^{-1}$ est une surjection $\llb 1, m \rrb \to \llb 1, n \rrb$.
            Pour chaque $i \in \llb 1, n \rrb$, on pose $r(i) =$ le plus petit élément de $q^{-1}(\{i\})$.
            $r$ est une injection $\llb 1, n \rrb \to \llb 1, m \rrb$.
            Donc $n \leq m$.
            Réciproquement, si $n \leq m$, alors on peut définir la surjection $\llb 1, m \rrb \to \llb 1, n \rrb$.
            $s(k) = \begin{cases} k, & \text{si} k \leq n \\ 1, & \text{si} k > n \end{cases}$
            Alors $u \circ s \circ v^{-1}$ est une surjection $F \to E$.
    \end{enumerate}
\end{proof}

\begin{cor}
    Soient $E, F$ deux ensembles finis,
    s'il existe une injection $E \to F$ et une injection $F \to E$,
    alors $|E| = |F|$.
\end{cor}

\begin{cor}
    Soit $E$ un ensemble fini.
    \begin{itemize}
        \item Si $A \subseteq E$, alors $A$ est fini et $|A| \leq |E|$.
        \item Si $f: E \to F$ est une application, alors $f(E)$ est fini et $|f(E)| \leq |E|$.
            De plus, $|f(E)|=|E|$ ssi $f$ est injective.
        \item Si $f: E \to E$, alors $f$ est injective $\iff f$ est surjective $\iff f$ est bijective.
    \end{itemize}
\end{cor}

\begin{proof}
    Soit $n=|E|$, alors on peut écrire $E=\{a_1, a_2, \dots, a_n\}$.
    \begin{itemize}
        \item Si $A \subseteq E$, alors $A$ s'écrit sous la forme
            $A=\{a_{i_1}, a_{i_2}, \dots, a_{i_p}\}$ avec $1 \le i_1 < i_2 < \dots < i_p \le n$.
            Donc $A$ est fini et $|A|=p \leq n$.
        \item On a $f(E)= \{f(a): a \in E\}=\{f(a_1), f(a_2), \dots, f(a_n)\}$,
            donc $f(E)$ contient au plus $n$ éléments, d'où $|f(E)| \leq n$.
            On a $|f(E)|= n$ ssi il n'y a pas de valeurs répétées parmi
            $f(a_1), f(a_2), \dots, f(a_n)$ càd $f$ est injective.
        \item Si $f$ est injective, selon le point 2, $|f(E)|=|E|$.
            Car $f(E) \subseteq E$ et $|f(E)|=|E|$, on a $f(E)=E$, donc $f$ est surjective.
            Si $f$ est surjective, alors $|f(E)|=|E|$, et selon le point 2, $f$ est injective.
            Donc $f$ est injective ssi $f$ est surjective.
    \end{itemize}
\end{proof}

\begin{rqq}
    L'hypothèse `` $E$ est fini '' est indispensable dans le corollaire 2.
    Si $E$ est infini, on a le contre-exemple : 
    \[
    f: \begin{array}{ccc}
    \mathbb{N} & \to & \mathbb{N} \\
    n & \mapsto & n+1
    \end{array}
    \]
    qui est injective mais non surjective.
\end{rqq}

\begin{cor}[Principe des tiroirs]
    Si on place $n$ objets dans $m$ tiroirs avec $n > m$, alors au moins un tiroir contient au moins 2 objets.
\end{cor}

\begin{proof}
    Soit $f: E = \text{l'ensemble des objets} \to F = \text{l'ensemble des tiroirs}$
    l'application qui décrit le placement des objets dans les tiroirs.
    On sait que $|E| > |F|$, alors $f$ n'est pas injective,
    càd $\exists y \in F, x_1, x_2 \in E$ tels que $x_1 \neq x_2$ mais $f(x_1)=f(x_2)=y$.
\end{proof}